\Problem{Caja de Ingenieros de Datos}

\Statement

En la Caja de Ingenieros de Datos se hace un sorteo para premiar el uso de las tarjetas de crédito. Las condiciones del premio son las siguientes:
\begin{itemize}
    \item Haber hecho almenos 3 compras.
    \item Que el total de las compras ascienda a 200€ o más.
    \item Que la última fecha de compra sea el 1 de Marzo de 2024 o posterior.
\end{itemize}

\medskip

Se necesita un programa que reciba a la entrada las compras hechas por todas las tarjetas de crédito de la Caja de Ingenieros de Datos y muestre a la salida los candidatos, es decir, aquellos números de tarjeta que cumplen con las condiciones para optar al premio.

\Input

Una secuencia, por orden cronológico, de compras con tarjeta de crédito que incluyen: a) la fecha, como un entero con el formato \verb|AAAAMMDD|; b) el número de tarjeta de crédito, y; c) el importe de la compra. Cada compra ocupa una línea con los 3 elementos separados por un espacio.

\Output

La salida debe ser una lista de los números de tarjeta candidatos al premio, por orden de menor a mayor, uno por línea.

\Sample

% \Observations


